\documentclass[11pt]{article}
\usepackage{amsmath, amssymb}
\usepackage{geometry}
\usepackage{parskip}
\usepackage{hyperref}
\geometry{margin=1in}

\title{ECON622 --- Project Notes}
\author{}
\date{}

\begin{document}
\maketitle

% ─────────────────────────────────────────────────────────────────────────────
\section{Setup}

\begin{itemize}
    \item $I$ = students, $J$ = course-instructor pairs, $c_j$ = TA slot capacity for course $j$
    \item Decision variables: $x_{ij} \in \{0,1\}$, where $x_{ij} = 1$ if student $i$ is assigned to course $j$
\end{itemize}

% ─────────────────────────────────────────────────────────────────────────────
\section{Data Generating Process}

Latent skill vectors $\theta_i \in \mathbb{R}^k$ (student) and $\phi_j \in \mathbb{R}^k$ (course), drawn i.i.d.\ from $\mathcal{N}(0, I_k)$.

Preference scores:
\[
    u_{ij} = \theta_i \cdot \phi_j + \varepsilon_{ij}, \qquad \varepsilon_{ij} \sim \mathcal{N}(0, \sigma^2)
\]
Course preferences are symmetric but use independent noise draws. Rankings are derived from scores.

% ─────────────────────────────────────────────────────────────────────────────
\section{MIP Formulation}

Represent the assignment as a binary matrix $X$ (students $\times$ courses).
The solver finds the $X$ that maximises a chosen objective subject to constraints.

\subsection*{Constraints}

\begin{align}
    \sum_j x_{ij} &\leq 1 \quad \forall\, i
        && \text{each student assigned at most once} \\[4pt]
    \sum_i x_{ij} &\leq c_j \quad \forall\, j
        && \text{course capacity} \\[4pt]
    x_{ij} &= 0 \quad \text{if forbidden}
        && \text{PhD constraint, rejection lists (encoded as a mask)}
\end{align}

\subsection*{Objective variants}

\begin{itemize}
    \item \textbf{Student utilitarian:} $\displaystyle\max \sum_{i,j} u_{ij}\, x_{ij}$
    \item \textbf{Course utilitarian:} $\displaystyle\max \sum_{i,j} v_{ij}\, x_{ij}$
    \item \textbf{Bilateral:} $\displaystyle\max \sum_{i,j} (u_{ij} + v_{ij})\, x_{ij}$
    \item \textbf{Egalitarian:} $\displaystyle\max\; t$ \quad (see below)
\end{itemize}

\subsection*{Egalitarian (min-max) formulation}

Introduce scalar $t$ representing the floor utility among matched students.
For each student $i$:
\[
    t \;\leq\; \sum_j u_{ij}\, x_{ij} \;+\; M\!\left(1 - \sum_j x_{ij}\right)
\]
where $M > \max_{i,j}|u_{ij}|$ is a big-M constant.

\begin{itemize}
    \item If student $i$ is matched ($\sum_j x_{ij} = 1$): constraint reduces to $t \leq \sum_j u_{ij} x_{ij}$. \textbf{Active.}
    \item If student $i$ is unmatched ($\sum_j x_{ij} = 0$): constraint becomes $t \leq \sum_j u_{ij} x_{ij} + M$, which is always satisfied. \textbf{Inactive.}
\end{itemize}

This prevents unmatched students from dragging $t$ to zero or below.

\subsection*{LP relaxation}

Relax $x_{ij} \in \{0,1\}$ to $x_{ij} \in [0,1]$ (continuous). Solve the LP, then round greedily: assign student $i$ to $\arg\max_j x_{ij}$ if that value exceeds $0.5$.

% ─────────────────────────────────────────────────────────────────────────────
\section{Algorithms Implemented}

\begin{enumerate}
    \item \textbf{Deferred Acceptance (student-proposing)} --- replicates the VSE approach.
          Many-to-one converted to one-to-one via allocation block expansion.
    \item \textbf{Serial Dictatorship} --- students pick in priority order; random order by default.
    \item \textbf{MIP / LP relaxation} --- four objective variants; Gurobi solver (HiGHS fallback).
\end{enumerate}

% ─────────────────────────────────────────────────────────────────────────────
\section{Evaluation Criteria}

\begin{itemize}
    \item Stability (no blocking pairs)
    \item Match rate
    \item Student rank distribution of matched course
    \item Welfare measures (total utility, min utility)
    \item Computational efficiency (solve time)
\end{itemize}

% ─────────────────────────────────────────────────────────────────────────────
\section{Thoughts / TODO}

\begin{itemize}
    \item[$\square$] MIP testing
    \item[$\square$] Comparison / evaluation framework across all algorithms
    \item[$\square$] Heterogeneous preferences experiment (vary $\sigma$, $k$)
    \item[$\square$] Visualisations
\end{itemize}

\end{document}
